%!TEX encoding = UTF8
%!TEX root = 0-notes.tex

\chapter*{Préface}


\section*{Syllabus}

6 ECTS. 30h CM, 30h TD.

\subsection*{Description}

Cette UE fait suite à l'UE de S1 (Analyse I) où on été introduits continuiét et dérivabilité des fonctions réelles, fonctions usuelles, et l'étude des suites réelles.
L'objectif est de poursuivre et d'approfondir le travail sur les suites et fonctions et d'introduire l'étude des séries numériques.

\subsection*{Objectifs}

 \begin{enumerate}
 	\item
 	Suites numériques.
 	\begin{enumerate}
 		\item
 		Relation de comparaison sur les suites (petit o, grand O, équivalent).
 		\item
 		Limite sup/limite inf, notion de valeur d'adhérence, suite de Cauchy (exemple de suite de Cauchy de rationnels qui ne converge pas dans $\Q$).
 		\item
 		Théorème de Bolzano-Weierstrass.
 		\item
 		Étude de suites récurrentes ($u_{n+1} = f(u_n)$).
	\end{enumerate}
 	\item
 	Fonctions réelles.
 	\begin{enumerate}
 		\item
 		Relation de comparaison (petit o, grand O, équivalent).
 		\item
 		Développements limites et formule de Taylor-Lagrange, Taylor-Young, développements limités usuels, opérations, applications des développements limités aux calculs de limites, intégralités usuelles, position relative d'une courbe par rapport à sa tangente, étude asymptotique.
 		\item
 		Régularité des fonctions : théorème des bornes atteintes, continuité uniforme, fonctions lipschitziennes, théorème de Heine.
	\end{enumerate}
 	\item
 	Étude des séries numériques.
 	\begin{enumerate}
 		\item
 		Séries géométriques et téléscopiques, cas simple avec calcul explicite des sommes partielles.
 		\item
 		Séries positives (relation de comparaison, séries de Riemann, critère de Cauchy/d'Alembert, critère de condensation, séries de Bertrand).
	\end{enumerate}
\end{enumerate}

\subsection*{Pré-requis nécessaires}

Programme de mathématiques du S1, et en particulier Analyse I, Logique, raisonnements, ensembles et applications, et Remédiation et Calculus.

\subsection*{Pré-requis recommandés}

Programme de mathématiques du S1.

\section*{Structure du cours}

Construction d'un cours de mathématiques : Exemple -- Remarque -- Définition -- Lemme -- Proposition -- Théorème -- Démonstration -- Corollaire.

Voir \emph{Mechanica, volume 1}, Leonhard Euler, 1736 (\href{https://scholarlycommons.pacific.edu/euler-works/15/}{E15}), pour un exemple de cours de mathématiques en latin.

Wolframalpha permet de faire du calcul exact.

\begin{definition*}
Associe un mot de vocabulaire à un objet mathématique.
Un mot de vocabulaire peut avoir un sens différent en mathématiques qu'en langage vernaculaire.
C'est par exemple le cas de « ensemble », « couple », « rationnel », « réel », ...
\end{definition*}

\begin{lemme*}
	Résultat préliminaire, parfois technique, utile à la démonstration d'une proposition plus importante.
\end{lemme*}

\begin{proposition*}
	Énoncé vrai et important.
\end{proposition*}

\begin{theorem*}
	Énoncé vrai et fondamental.
\end{theorem*}

\begin{corollaire*}
	Résultat particulier, découlant souvent quasi directement d'un théorème ou d'une proposition.
\end{corollaire*}

%\newcounter{preface}
\begin{Exercise}[label=1]%[counter=preface]
	Exemple d'exercice d'application directe.
\end{Exercise}
\begin{Exercise}[difficulty=1, label=2]%, counter=preface]
	Exemple de problème.
\end{Exercise}
\begin{Exercise}[difficulty=2, label=3]%, counter=preface]
	Exemple de problème d'approfondissement.
\end{Exercise}
\shipoutExercise
\begin{Answer}[ref=1]
	Solution 1.
\end{Answer}
\begin{Answer}[ref=2]
	Solution 2.
\end{Answer}
\begin{Answer}[ref=3]
	Solution 3.
\end{Answer}
\begin{multicols}{3}
\shipoutAnswer
\end{multicols}
